\chapter{Sintaxis \emph{graph}: el separador \texttt{;}}
\label{cap:dsl-graph}

\begin{intuicion}
La sintaxis chain (capítulo anterior) sólo expresa cadenas
lineales. Para hipótesis tipo estrella o árbol ---``un usuario
que se autentica en tres hosts distintos''--- la cadena no alcanza.
La extensión L1 stage 3c agregó un solo separador, \texttt{;}, que
abre las puertas a hipótesis no-cadena. Una hipótesis con
\texttt{;} en cualquier parte se marca como \emph{shape graph}; el
parser construye un hipergrafo, el planner detecta
\(\alpha\)-aciclidad, y el matcher dispatcha a Yannakakis tree-walk
(capítulo~\ref{cap:yannakakis-tree}). Este capítulo documenta la
sintaxis y el flujo del parser.
\end{intuicion}

\section{El cambio mínimo: agregar \texttt{;}}

\begin{lstlisting}
User c -[Auth]-> Host h1; c -[Auth]-> Host h2; c -[Auth]-> Host h3
\end{lstlisting}

\noindent
Esa es la hipótesis ``un usuario que se autentica en al menos tres
hosts distintos'' (la \emph{star query}). Tres cláusulas separadas
por \texttt{;}; las tres comparten el binding \texttt{c}.

Compará con la equivalente en chain pura:
\begin{lstlisting}
User c -[Auth]-> Host  HAVING count(distinct last) >= 3
\end{lstlisting}

\noindent
La versión chain es más concisa pero menos expresiva: enumera
\emph{cualquier} count de hosts \(\geq 3\), sin nombrarlos. La
graph version los nombra (\(h_1, h_2, h_3\)) y permite predicados
distintos sobre cada uno.

\section{Reglas de las cláusulas tras \texttt{;}}

\subsection{Bare binding reference}

\begin{lstlisting}
User c -[Auth]-> Host h1; c -[Auth]-> Host h2
\end{lstlisting}

\noindent
La segunda cláusula arranca con \texttt{c}, que NO es un type sino
un binding ya declarado. El parser:

\begin{enumerate}
    \item Peeka el primer token de la cláusula post-\texttt{;}.
    \item Si matchea un binding declarado en cláusulas previas
    (\texttt{c} sí), retorna \texttt{(existing\_type, Some(name))}
    al loop principal.
    \item El loop arranca el chain step desde ese vértice
    compartido: hereda el \texttt{EntityType} (\texttt{User}) del
    binding original.
\end{enumerate}

\subsection{Fresh type declaration}

\begin{lstlisting}
User c -[Auth]-> Host h1; Process p -[Spawn]-> File f
\end{lstlisting}

\noindent
La segunda cláusula arranca con \texttt{Process p}, que es un nuevo
type declaration. El parser introduce un vértice fresco
independiente del primero. La hipótesis representa: ``ambos
patrones existen en el grafo, sin conexión entre ellos''.

\begin{nota}
Este caso (cláusulas \emph{disconectadas} unidas sólo por
\texttt{;}) es legal pero raro: el matcher las trata como un
\emph{producto cartesiano} de matches independientes. Output:
$|\text{matches del primero}| \times |\text{matches del segundo}|$.
Útil sólo cuando el analista quiere correlacionar patrones
distintos.
\end{nota}

\subsection{Disambiguación}

¿Qué hace el parser si el primer token de la cláusula post-\texttt{;}
es ambiguo? Reglas:

\begin{enumerate}
    \item Si matchea un \emph{builtin EntityType}
    (\texttt{User}, \texttt{Host}, \texttt{IP}, etc.) $\Rightarrow$
    siempre se parsea como type declaration. Sombrea cualquier
    binding del mismo nombre. (Razón: los EntityTypes son
    palabras reservadas; usar ``User'' como binding sería
    confuso.)
    \item Si matchea un binding ya declarado $\Rightarrow$ bare
    binding ref.
    \item Si no matchea ninguno $\Rightarrow$ \textbf{error} con
    \emph{hint}: ``binding desconocido; los declarados son: \dots''.
\end{enumerate}

\noindent
El error con hint es deliberado: descarta el fallback silencioso a
\texttt{Other} (capítulo~\ref{cap:entity-relation}). Los typos del
analista (\texttt{Hsot} en vez de \texttt{Host}) tienen que
saltar como error, no como ``vértice de tipo Other''.

\section{Estructura interna del parser}

\begin{figure}[H]
\centering
\begin{tikzpicture}[
    node distance=0.5cm,
    every node/.style={font=\scriptsize, align=center},
    state/.style={draw=accent, thick, rounded corners=1mm,
                  minimum height=0.7cm, minimum width=2.6cm,
                  fill=accentlight}
]
\node[state] (start) {\texttt{parse()}};
\node[state, below=of start] (head) {\texttt{parse\_entity\_type\_with\_binding}};
\node[state, below=of head] (chain) {\texttt{parse\_chain\_steps}};
\node[state, below=of chain] (sep) {peek \texttt{;}};
\node[state, below left=of sep] (clause) {\texttt{parse\_clause\_head}};
\node[state, below right=of sep] (trail) {\texttt{parse\_trailing\_clauses}};

\draw[->, >=Latex, accent, thick] (start) -- (head);
\draw[->, >=Latex, accent, thick] (head) -- (chain);
\draw[->, >=Latex, accent, thick] (chain) -- (sep);
\draw[->, >=Latex, accent, thick] (sep) -- node[left]{si \texttt{;}} (clause);
\draw[->, >=Latex, accent, thick, dashed]
    (clause.east) to[out=0,in=180] (chain.west);
\draw[->, >=Latex, accent, thick] (sep) -- node[right]{else} (trail);
\end{tikzpicture}
\caption{Flujo del parser. \texttt{parse\_chain\_steps} es el bucle
interior reusable; cada \texttt{;} dispara un re-entry vía
\texttt{parse\_clause\_head}, que decide entre bare binding ref y
type declaration.}
\label{fig:dsl-parser-flow}
\end{figure}

\subsection{Cómo se lee cada nodo}

\begin{itemize}
    \item \texttt{parse\_entity\_type\_with\_binding}
    (\emph{first clause head}): consume un \texttt{Type}
    obligatorio y, si lo sigue un identificador no reservado, lo
    registra como binding. Es la única forma de declarar el primer
    vértice ---las hipótesis no pueden empezar con un bare binding
    ref porque no hay nada declarado todavía---.

    \item \texttt{parse\_chain\_steps}: bucle que consume \(N\)
    pasos \texttt{-[rel]-> Type binding} encadenados. Termina cuando
    ve EOF, \texttt{\{} (suffix \texttt{\{k=N\}} o predicate block),
    \texttt{;} (clause separator), o un identificador alfabético
    (HAVING/WITHIN). Cada terminación delega al manejador
    correspondiente.

    \item \texttt{parse\_clause\_head} (\emph{tras un \texttt{;}}):
    peeka el identificador líder y decide:

    \begin{itemize}
        \item Si matchea un builtin type name (\texttt{User},
        \texttt{Host}, \texttt{IP}, etc.) $\Rightarrow$ rewind y
        delega a \texttt{parse\_entity\_type\_with\_binding}.
        \item Si matchea un binding ya declarado $\Rightarrow$
        retorna \texttt{(existing\_type, Some(name))} para que el
        bucle de \texttt{parse\_chain\_steps} arranque desde ese
        vértice compartido.
        \item Si no es ninguno de los dos $\Rightarrow$ error con
        hint listando los bindings válidos.
    \end{itemize}

    \item \texttt{parse\_trailing\_clauses}: consume \emph{HAVING}
    y \emph{WITHIN} en cualquier orden hasta encontrar
    \texttt{\{k=N\}} o EOF. Estos no pueden re-entrar el bucle de
    chain steps; son estrictamente post-fix a la última cláusula
    del cuerpo.
\end{itemize}

\section{Ejemplos}

\begin{ejemplo}[Estrella de tres brazos]
\begin{lstlisting}
User c -[Auth]-> Host h1; c -[Auth]-> Host h2; c -[Auth]-> Host h3
\end{lstlisting}

Tres cláusulas, todas comparten \(c\). Shape graph $\Rightarrow$
Yannakakis tree-walk. El join tree es una estrella: nodo central
es step 0; los otros dos son hijos.
\end{ejemplo}

\begin{ejemplo}[Triángulo de relaciones]
\begin{lstlisting}
User a -[Knows]-> User b; b -[Knows]-> User c; c -[Knows]-> a
\end{lstlisting}

Tres cláusulas formando ciclo: \(a \to b \to c \to a\). El
hipergrafo es cíclico (no \(\alpha\)-acíclico). El planner detecta
y cae al DFS legacy ---no es \(K_4\), no aplica LFTJ; no es
\(\alpha\)-acíclico, no aplica Yannakakis.
\end{ejemplo}

\begin{ejemplo}[Árbol de cuatro niveles]
\begin{lstlisting}
User u -[Auth]-> Host h; h -[Hosts]-> Process p; p -[Spawn]-> Process c
\end{lstlisting}

Tres cláusulas conectadas en árbol (\(u \to h \to p \to c\)). Es
\(\alpha\)-acíclico $\Rightarrow$ Yannakakis. Notar que esta
hipótesis se podría haber escrito como chain pura
(\texttt{User -[Auth]-> Host -[Hosts]-> Process -[Spawn]-> Process});
la versión graph no aporta nada nuevo. La forma chain es más
concisa.

(El motor detecta esto: si la graph hypothesis es estrictamente
una chain disfrazada, podría re-escribirla como chain. Hoy no lo
hace; ambas formas pasan por sus respectivos matchers, y los
benches no muestran diferencia significativa.)
\end{ejemplo}

\section{Limitaciones}

\begin{itemize}
    \item \textbf{No hay disjunción intra-cláusula}. Cada cláusula
    es una chain pura. Si querés ``\(c\) tiene un Auth A o B'',
    necesitás dos hipótesis separadas.
    \item \textbf{No hay subqueries}. Cada cláusula es plana; no
    hay anidamiento.
    \item \textbf{Sólo \(\alpha\)-acíclicos eficientes}. Si la
    hipótesis introduce un ciclo (como el ejemplo del triángulo),
    el motor cae al DFS, perdiendo el lift de Yannakakis.
\end{itemize}

\section{Resumen}

\begin{itemize}
    \item Sintaxis graph: chains separadas por \texttt{;}, todas
    comparten variables (bindings) cuando se referencian.
    \item Cada cláusula post-\texttt{;} arranca con type declaration
    (vértice fresco) o bare binding ref (vértice compartido).
    \item Errores explícitos en typos: ``binding desconocido''
    en vez de fallback silencioso.
    \item \texttt{shape == Graph} si y sólo si hay \texttt{;}.
    Yannakakis tree-walk dispatcha si \(\alpha\)-acíclico; DFS si
    cíclico.
\end{itemize}

\begin{ejercicios}
\begin{enumerate}
    \item \starred{} Escribí en sintaxis graph: ``un usuario que
    crea un proceso, y ese proceso escribe un archivo, y ese
    proceso también lanza otro proceso''. Identificá los bindings
    compartidos.

    \item Para la hipótesis del ejercicio 1, dibujá el join tree.
    ¿Cuál es la root?

    \item \starred{} ¿La hipótesis
    \texttt{User c -[Auth]-> Host h; c -[Auth]-> Host h}
    (dos cláusulas idénticas) es legal? Si sí, ¿qué encuentra el
    matcher?

    \item Diseñá una hipótesis graph cíclica de tres pasos.
    ¿Cómo la detecta el planner? ¿A qué operador cae?

    \item Si una hipótesis tiene 5 cláusulas separadas por
    \texttt{;}, todas \emph{disconectadas} (sin bindings
    compartidos), ¿qué tamaño de output puede tener? Cota teórica.
\end{enumerate}
\end{ejercicios}
